% =============================================================================
% INTRODUCTION (social-science structure: Hook → Gap → RQ → Contribution → Findings → Roadmap)
% Follows introduction reference of social-science-paper skill.
% =============================================================================

\section{Introduction}
\label{sec:introduction}

Many studies of enterprise digitalization, environmental performance, and productivity rely on revenue-based or accounting proxies when direct measures are unavailable: digitalization is proxied by performance rankings or composite indices, total factor productivity (TFP) by accounting returns such as ROE, and environmental performance by growth outcomes or ESG scores. When regressor and outcome share the same underlying variable (e.g., revenue), \emph{same-source bias} arises; together with construct mismatch, it makes the resulting coefficients uninterpretable for causal inference. The regressions then capture \emph{revenue dynamics} rather than the target constructs. Whether digitalization improves both TFP and environmental performance in heavily polluting firms (a ``double dividend'') is policy-relevant and under-researched; yet the prevalence and formal structure of this measurement problem have received little attention.

Research in this area often turns to proxies because text-based digitalization indices, production-function TFP (e.g., Levinsohn--Petrin), and emission or green TFP data are frequently missing or costly, especially in emerging markets and for heavily polluting firms. What remains unclear is \emph{why} such proxies mislead, \emph{how widespread} the problem is, and whether researchers can diagnose same-source proxy bias in a principled way. Without a formal characterization of identification failure and a usable diagnostic, the literature cannot tell whether estimated associations reflect the target constructs or merely revenue and accounting outcomes.

\textbf{Our research question is: Why do revenue-based proxies mislead (or fail to identify) the relationship between digitalization and the double dividend, and how widespread is this measurement problem?} This paper does \emph{not} estimate that substantive relationship; it addresses the methodological gap by documenting prevalence, formalizing same-source proxy bias, and providing a diagnostic that researchers can apply to their own designs.

We make four contributions. First, we \emph{document prevalence} through a structured literature audit of 100 firm-level studies: 83\% use at least one revenue-related or accounting-based proxy; we provide a replicable search protocol and journal-tier classification (Section~\ref{sec:related}, \ref{app:audit-protocol}). Second, we \emph{formalize} same-source proxy bias: Proposition~1 establishes identification failure when regressor and outcome are functions of the same source; Corollary~1 shows that under monotone proxies the bias has a systematic sign (so a negative coefficient need not imply a negative causal effect) and is non-vanishing; Theorem~2 shows that the $t$-statistic can diverge with sample size, so proxy-based false positives become more convincing in larger samples. We also provide a \emph{three-step diagnostic} with a practical checklist (Section~\ref{sec:formal-diagnostic}). Third, we \emph{demonstrate} measurement failure with a transparent proxy-based illustration using open data on Chinese A-share listed firms in heavily polluting industries (empirical illustration only) and a controlled experiment that implements the diagnostic. Fourth, we offer a \emph{blueprint} for valid inference with direct measures. Taken together, we show that same-source proxies not only fail to identify the target effect but can produce directionally systematic, asymptotically significant coefficients when the true effect is zero, so large-sample significance is no evidence against measurement failure and can be a symptom of it.

The audit confirms that proxy use is widespread. The proxy illustration shows that regressions with revenue rank, ROE, and revenue growth capture revenue-based variation rather than digitalization effects. The controlled experiment shows that when true digitalization is by design uncorrelated with revenue, the same proxy regression still yields a statistically significant coefficient, demonstrating that same-source structure alone can generate spurious evidence.

The remainder of this paper proceeds as follows. Section~\ref{sec:related} reviews the measurement gap and presents the literature audit. Section~\ref{sec:theory} clarifies target versus implemented constructs. Section~\ref{sec:methods} sets out the target and implemented designs, the formal characterization and diagnostic, and the sample construction. Section~\ref{sec:results} presents the proxy illustration and the controlled experiment. Section~\ref{sec:discussion} and Section~\ref{sec:conclusion} discuss implications and conclude. Additional regression tables (mediation, moderation, robustness, weak-instrument IV) are in the appendix.
